% !TEX spellcheck = en_US



% ~~~~~~~~~~~~~~~~~~~~~~~~~~~~~~~~~~~~~~~ V
% (NC) 2026 Haluk Bingol
% github.com/halukbingol/LaTeX-Templates
%
% Licensees may copy, distribute, display, and perform the work and make derivative works 
% and remixes based on it only for non-commercial purposes. 
% ~~~~~~~~~~~~~~~~~~~~~~~~~~~~~~~~~~~~~~~ A




% =====================================================================
%  Trees -- Sophomore Level: Generic Trees, Iterative Traversals and Search Trees
%  ---------------------------------------------------------------------
%  Built on hblib/Template/hblibTemplate.tex with the libraries in ../../hblib/.
%  Programs and their real outputs are in codeoutput/:
%     codeoutput/<Name>.java  a program (classes without main() are only compiled)
%     codeoutput/<Name>.sh    a terminal session, run in a copy of codeoutput/<Name>/
%     codeoutput/<Name>.txt   what it printed, made by:  bash run_examples.sh
% =====================================================================

\documentclass[aspectratio=169,10pt,compress]{beamer}

% ~~~~~~~~~~~~~~~~~~~~~~~~~~~~~~~~~~~~~~~~~~~~~~~~~~~~~~~~~~~~~~~~~ V hblib
	\usepackage{../../hblib/hblibPresentation} % lib presentation: theme, i/N, \hSection, algorithm2e, ...
	\usepackage{../../hblib/hblibCodeoutput} % lib code + output: \lstcode, \lstterminal, ...
	\usepackage{../../hblib/hblibDatastructures} % lib data structures: \hString, \hStack, \hQueue
% ~~~~~~~~~~~~~~~~~~~~~~~~~~~~~~~~~~~~~~~~~~~~~~~~~~~~~~~~~~~~~~~~~ A hblib


% xcolor, array (and the L{width} column), listings, tikz and algorithm2e come from the hblib libraries
\usepackage[T1]{fontenc}
\usepackage{lmodern}
\usepackage{booktabs}

\newcommand{\win}{\textbf{Windows}}
\newcommand{\mac}{\textbf{macOS}}

% algorithm2e set-up (the package is loaded by hblibPresentation)
\SetAlCapFnt{\small}
\SetAlCapNameFnt{\small}
\SetKwProg{Fn}{function}{}{end}
\SetKw{KwError}{error}

% ---------------------------------------------------------------------
\title{Trees}
\subtitle{Sophomore Level: Generic Trees, Iterative Traversals and Search Trees}
\author[Bingol]{Haluk O. Bingol}

\institute[FCIS]{
    Faculty of Computer and Information Sciences\\
    Yeditepe University
}
\date{
	{\tiny \hbTimeStamp}
}

\begin{document}




% ~~~~~~~~~~~~~~~~~~~~~~~~~~~~~~~~~~~~~~ V frm
\begin{frame}[t,fragile]
  \titlepage
\end{frame}
% ~~~~~~~~~~~~~~~~~~~~~~~~~~~~~~~~~~~~~~ A frm




% ~~~~~~~~~~~~~~~~~~~~~~~~~~~~~~~~~~~~~~ V frm
\begin{frame}[t,fragile] \frametitle{Outline}
  \begin{columns}[T]
    \begin{column}{0.48\textwidth}
      \tableofcontents[sections={1-4}, hideallsubsections]
    \end{column}
    \begin{column}{0.48\textwidth}
      \tableofcontents[sections={5-8}, hideallsubsections]
    \end{column}
  \end{columns}
\end{frame}
% ~~~~~~~~~~~~~~~~~~~~~~~~~~~~~~~~~~~~~~ A frm

% ~~~~~~~~~~~~~~~~~~~~~~~~~~~~~~~~~~~~~~ V
\hSection[Generic]{A Generic Binary Tree}
% ~~~~~~~~~~~~~~~~~~~~~~~~~~~~~~~~~~~~~~ A




% ~~~~~~~~~~~~~~~~~~~~~~~~~~~~~~~~~~~~~~ V frm
\begin{frame}[t,fragile] \frametitle{\texttt{Node<T>}: a tree of anything}
  \begin{columns}[T]
    \begin{column}{0.6\textwidth}
      \lstcode[basicstyle=\ttfamily\tiny]{codeoutput/Node.java}{Java}
    \end{column}
    \begin{column}{0.38\textwidth}
      \begin{center}
      \begin{tikzpicture}[htree, level distance=0.95cm, level 1/.style={sibling distance=2.0cm}, level 2/.style={sibling distance=1.0cm}, level 3/.style={sibling distance=0.55cm}, scale=0.9, transform shape]
  \node[hnode] {A} child {node[hnode] {B} child {node[hnode] {D}} child {node[hnode] {E}}} child {node[hnode] {C} child[missing] child {node[hnode] {F} child {node[hnode] {G}} child[missing]}};
\end{tikzpicture}
      \end{center}
      \begin{itemize}\small
        \item \texttt{T}: \texttt{String}, \texttt{Integer}, \texttt{Student}, \dots
        \item \texttt{leaf} is a \textbf{generic method}: \texttt{<T>} before the return type.
      \end{itemize}
    \end{column}
  \end{columns}
\end{frame}
% ~~~~~~~~~~~~~~~~~~~~~~~~~~~~~~~~~~~~~~ A frm




% ~~~~~~~~~~~~~~~~~~~~~~~~~~~~~~~~~~~~~~ V frm
\begin{frame}[t,fragile] \frametitle{Traversals with a visitor}
  \begin{columns}[T]
    \begin{column}{0.49\textwidth}
      \lstcode[basicstyle=\ttfamily\tiny, lastline=26]{codeoutput/Walk.java}{Java}
    \end{column}
    \begin{column}{0.49\textwidth}
      \lstcode[basicstyle=\ttfamily\tiny, firstline=27, firstnumber=27]{codeoutput/Walk.java}{Java}
    \end{column}
  \end{columns}

  \small \texttt{Consumer<T> visit}: the traversal decides the \textbf{order}, the caller decides \textbf{what to do} with each value.
\end{frame}
% ~~~~~~~~~~~~~~~~~~~~~~~~~~~~~~~~~~~~~~ A frm




% ~~~~~~~~~~~~~~~~~~~~~~~~~~~~~~~~~~~~~~ V frm
\begin{frame}[t,fragile] \frametitle{Passing actions as lambdas}
  \lstcoderightstep[0.64][basicstyle=\ttfamily\tiny][basicstyle=\ttfamily\tiny]{codeoutput/UseWalk.java}{Java}{codeoutput/UseWalk.txt}

  \vspace{0.2em}
  \small \texttt{out::add}, \texttt{v -> count[0]++}: one traversal, many uses. (A lambda cannot change a local \texttt{int}, so \texttt{count} is an array.)
\end{frame}
% ~~~~~~~~~~~~~~~~~~~~~~~~~~~~~~~~~~~~~~ A frm

% ~~~~~~~~~~~~~~~~~~~~~~~~~~~~~~~~~~~~~~ V
\hSection[Iterative]{Traversals without Recursion}
% ~~~~~~~~~~~~~~~~~~~~~~~~~~~~~~~~~~~~~~ A




% ~~~~~~~~~~~~~~~~~~~~~~~~~~~~~~~~~~~~~~ V frm
\begin{frame}[t,fragile] \frametitle{Preorder and inorder with a stack}
  \begin{columns}[T]
    \begin{column}{0.49\textwidth}
      \lstcode[basicstyle=\ttfamily\tiny, lastline=16]{codeoutput/Iterative.java}{Java}
    \end{column}
    \begin{column}{0.49\textwidth}
      \lstcode[basicstyle=\ttfamily\tiny, firstline=17, firstnumber=17, lastline=29]{codeoutput/Iterative.java}{Java}
    \end{column}
  \end{columns}
\end{frame}
% ~~~~~~~~~~~~~~~~~~~~~~~~~~~~~~~~~~~~~~ A frm




% ~~~~~~~~~~~~~~~~~~~~~~~~~~~~~~~~~~~~~~ V frm
\begin{frame}[t,fragile] \frametitle{Levels with a queue}
  \begin{columns}[T]
    \begin{column}{0.49\textwidth}
      \lstcode[basicstyle=\ttfamily\tiny, firstline=30, firstnumber=30, lastline=46]{codeoutput/Iterative.java}{Java}
    \end{column}
    \begin{column}{0.49\textwidth}
      \lstcode[basicstyle=\ttfamily\tiny, firstline=47, firstnumber=47]{codeoutput/Iterative.java}{Java}
    \end{column}
  \end{columns}
\end{frame}
% ~~~~~~~~~~~~~~~~~~~~~~~~~~~~~~~~~~~~~~ A frm




% ~~~~~~~~~~~~~~~~~~~~~~~~~~~~~~~~~~~~~~ V frm
\begin{frame}[t,fragile] \frametitle{How the iterative versions work}
  \begin{columns}[T]
    \begin{column}{0.5\textwidth}
      \lstoutput[basicstyle=\ttfamily\tiny]{codeoutput/Iterative.txt}
      \begin{itemize}\small
        \item \textbf{preorder}: pop a node, visit it, push right then left: left is popped first.
        \item \textbf{inorder}: go left as far as possible pushing nodes; pop, visit, then continue with its right subtree.
        \item \textbf{levels}: the queue's size at the start of a level is the number of nodes in that level.
      \end{itemize}
    \end{column}
    \begin{column}{0.46\textwidth}
      \small Stack contents during \textbf{iterative inorder} (top on the right):
      \begin{center}\scriptsize
      \begin{tabular}{lll}
        \toprule
        \textbf{step}             & \textbf{stack}    & \textbf{output} \\
        \midrule
        push A, B, D              & A B D             & \\
        pop D                     & A B               & D \\
        pop B, push E             & A E               & D B \\
        pop E                     & A                 & D B E \\
        pop A, push C             & C                 & \dots{} A \\
        pop C, push F, G          & F G               & \dots{} C \\
        pop G, pop F              &                   & \dots{} G F \\
        \bottomrule
      \end{tabular}
      \end{center}
      \small No recursion: the depth is limited by the heap, not by the call stack.
    \end{column}
  \end{columns}
\end{frame}
% ~~~~~~~~~~~~~~~~~~~~~~~~~~~~~~~~~~~~~~ A frm




% ~~~~~~~~~~~~~~~~~~~~~~~~~~~~~~~~~~~~~~ V frm
\begin{frame}[t,fragile] \frametitle{An iterator: inorder, one value at a time}
  \begin{columns}[T]
    \begin{column}{0.6\textwidth}
      \lstcode[basicstyle=\ttfamily\tiny]{codeoutput/InorderIterator.java}{Java}
    \end{column}
    \begin{column}{0.38\textwidth}
      \begin{itemize}\small
        \item The iterative inorder, cut into pieces: \texttt{next()} does one step.
        \item The stack holds at most $h + 1$ nodes: $O(h)$ memory, no list of all values.
        \item Each node is pushed and popped once: all of \texttt{next()} together are $O(n)$, so $O(1)$ \textbf{amortized} per call.
        \item With \texttt{Iterable}, a tree works in a for-each loop.
      \end{itemize}
    \end{column}
  \end{columns}
\end{frame}
% ~~~~~~~~~~~~~~~~~~~~~~~~~~~~~~~~~~~~~~ A frm

% ~~~~~~~~~~~~~~~~~~~~~~~~~~~~~~~~~~~~~~ V
\hSection[BST]{Binary Search Trees}
% ~~~~~~~~~~~~~~~~~~~~~~~~~~~~~~~~~~~~~~ A




% ~~~~~~~~~~~~~~~~~~~~~~~~~~~~~~~~~~~~~~ V frm
\begin{frame}[t,fragile] \frametitle{\texttt{Bst<K>}: insert and search}
  \begin{columns}[T]
    \begin{column}{0.49\textwidth}
      \lstcode[basicstyle=\ttfamily\tiny, lastline=14]{codeoutput/Bst.java}{Java}
    \end{column}
    \begin{column}{0.49\textwidth}
      \lstcode[basicstyle=\ttfamily\tiny, firstline=15, firstnumber=15, lastline=24]{codeoutput/Bst.java}{Java}
    \end{column}
  \end{columns}
\end{frame}
% ~~~~~~~~~~~~~~~~~~~~~~~~~~~~~~~~~~~~~~ A frm




% ~~~~~~~~~~~~~~~~~~~~~~~~~~~~~~~~~~~~~~ V frm
\begin{frame}[t,fragile] \frametitle{\texttt{Bst<K>}: delete, iterate, print}
  \begin{columns}[T]
    \begin{column}{0.49\textwidth}
      \lstcode[basicstyle=\ttfamily\tiny, firstline=25, firstnumber=25, lastline=42]{codeoutput/Bst.java}{Java}
    \end{column}
    \begin{column}{0.49\textwidth}
      \lstcode[basicstyle=\ttfamily\tiny, firstline=43, firstnumber=43]{codeoutput/Bst.java}{Java}
    \end{column}
  \end{columns}
\end{frame}
% ~~~~~~~~~~~~~~~~~~~~~~~~~~~~~~~~~~~~~~ A frm




% ~~~~~~~~~~~~~~~~~~~~~~~~~~~~~~~~~~~~~~ V frm
\begin{frame}[t,fragile] \frametitle{Deleting a node with two children}
  \begin{columns}[T]
    \begin{column}{0.32\textwidth}
      \centering
      \begin{tikzpicture}[htree, level distance=0.95cm, level 1/.style={sibling distance=1.9cm}, level 2/.style={sibling distance=1.0cm}, level 3/.style={sibling distance=0.55cm}, scale=0.8, transform shape]
  \node[hnode, hhot] {50} child {node[hnode] {30} child {node[hnode] {20}} child {node[hnode] {40} child {node[hnode] {35}} child[missing]}} child {node[hnode] {70} child {node[hnode, hdone] {60}} child {node[hnode] {80}}};
\end{tikzpicture}\\
      \small before: delete \texttt{50}
    \end{column}
    \begin{column}{0.32\textwidth}
      \centering
      \begin{tikzpicture}[htree, level distance=0.95cm, level 1/.style={sibling distance=1.9cm}, level 2/.style={sibling distance=1.0cm}, level 3/.style={sibling distance=0.55cm}, scale=0.8, transform shape]
  \node[hnode, hhot] {60} child {node[hnode] {30} child {node[hnode] {20}} child {node[hnode] {40} child {node[hnode] {35}} child[missing]}} child {node[hnode] {70} child[missing] child {node[hnode] {80}}};
\end{tikzpicture}\\
      \small after: \texttt{60} moved up
    \end{column}
    \begin{column}{0.32\textwidth}
      \begin{enumerate}\small
        \item no child: remove the node
        \item one child: the child takes its place
        \item two children: copy the \textbf{smallest key of the right subtree} (the inorder successor, red) into the node, then delete it there; it has no left child, so case 1 or 2
      \end{enumerate}
    \end{column}
  \end{columns}
\end{frame}
% ~~~~~~~~~~~~~~~~~~~~~~~~~~~~~~~~~~~~~~ A frm




% ~~~~~~~~~~~~~~~~~~~~~~~~~~~~~~~~~~~~~~ V frm
\begin{frame}[t,fragile] \frametitle{Using the search tree}
  \lstcoderightstep[0.6][basicstyle=\ttfamily\tiny][basicstyle=\ttfamily\tiny]{codeoutput/UseBst.java}{Java}{codeoutput/UseBst.txt}

  \vspace{0.2em}
  \small All operations follow one path from the root: $O(h)$. For random keys the average depth is about $1.39 \log_{2} n$; for sorted keys $h = n - 1$ (junior deck).
\end{frame}
% ~~~~~~~~~~~~~~~~~~~~~~~~~~~~~~~~~~~~~~ A frm




% ~~~~~~~~~~~~~~~~~~~~~~~~~~~~~~~~~~~~~~ V frm
\begin{frame}[t,fragile] \frametitle{Checking the BST property: a classic bug}
  \begin{columns}[T]
    \begin{column}{0.49\textwidth}
      \lstcode[basicstyle=\ttfamily\tiny, lastline=14]{codeoutput/CheckBst.java}{Java}
    \end{column}
    \begin{column}{0.49\textwidth}
      \lstcode[basicstyle=\ttfamily\tiny, firstline=15, firstnumber=15]{codeoutput/CheckBst.java}{Java}
    \end{column}
  \end{columns}
\end{frame}
% ~~~~~~~~~~~~~~~~~~~~~~~~~~~~~~~~~~~~~~ A frm




% ~~~~~~~~~~~~~~~~~~~~~~~~~~~~~~~~~~~~~~ V frm
\begin{frame}[t,fragile] \frametitle{Checking the BST property: the output}
  \begin{columns}[T]
    \begin{column}{0.3\textwidth}
      \begin{center}
      \begin{tikzpicture}[htree, level distance=0.95cm, level 1/.style={sibling distance=1.8cm}, level 2/.style={sibling distance=1.0cm}, scale=1.0, transform shape]
  \node[hnode] {50} child {node[hnode] {30} child {node[hnode] {20}} child {node[hnode, hhot] {60}}} child {node[hnode] {70}};
\end{tikzpicture}
      \end{center}
    \end{column}
    \begin{column}{0.66\textwidth}
      \lstoutput[basicstyle=\ttfamily\scriptsize]{codeoutput/CheckBst.txt}
      \begin{itemize}\small
        \item Every parent--child pair is in order, but \texttt{60} is in the \textbf{left} subtree of \texttt{50}.
        \item The BST rule is about \textbf{all} nodes of a subtree: pass the allowed range $(lo, hi)$ down.
        \item Equivalent test: the inorder sequence must be strictly increasing.
      \end{itemize}
    \end{column}
  \end{columns}
\end{frame}
% ~~~~~~~~~~~~~~~~~~~~~~~~~~~~~~~~~~~~~~ A frm

% ~~~~~~~~~~~~~~~~~~~~~~~~~~~~~~~~~~~~~~ V
\hSection[Rebuild]{Rebuilding and Saving Trees}
% ~~~~~~~~~~~~~~~~~~~~~~~~~~~~~~~~~~~~~~ A




% ~~~~~~~~~~~~~~~~~~~~~~~~~~~~~~~~~~~~~~ V frm
\begin{frame}[t,fragile] \frametitle{Preorder + inorder determine the tree}
  \lstcoderightstep[0.62][basicstyle=\ttfamily\tiny][basicstyle=\ttfamily\tiny]{codeoutput/Rebuild.java}{Java}{codeoutput/Rebuild.txt}

  \vspace{0.2em}
  \small The first preorder value is the root; its position in inorder splits the rest into left and right. (Preorder alone, or pre- and postorder, are not enough.)
\end{frame}
% ~~~~~~~~~~~~~~~~~~~~~~~~~~~~~~~~~~~~~~ A frm




% ~~~~~~~~~~~~~~~~~~~~~~~~~~~~~~~~~~~~~~ V frm
\begin{frame}[t,fragile] \frametitle{Writing a tree to text and reading it back}
  \lstcoderightstep[0.6][basicstyle=\ttfamily\tiny][basicstyle=\ttfamily\tiny]{codeoutput/Serialize.java}{Java}{codeoutput/Serialize.txt}

  \vspace{0.2em}
  \small Preorder \textbf{with markers for \texttt{null}} is enough on its own: the \texttt{\#} tell the reader where each subtree ends.
\end{frame}
% ~~~~~~~~~~~~~~~~~~~~~~~~~~~~~~~~~~~~~~ A frm

% ~~~~~~~~~~~~~~~~~~~~~~~~~~~~~~~~~~~~~~ V
\hSection[N-ary]{Generic N-ary Trees}
% ~~~~~~~~~~~~~~~~~~~~~~~~~~~~~~~~~~~~~~ A




% ~~~~~~~~~~~~~~~~~~~~~~~~~~~~~~~~~~~~~~ V frm
\begin{frame}[t,fragile] \frametitle{\texttt{TreeNode<T>} with a list of children}
  \begin{columns}[T]
    \begin{column}{0.49\textwidth}
      \lstcode[basicstyle=\ttfamily\tiny, lastline=18]{codeoutput/TreeNode.java}{Java}
    \end{column}
    \begin{column}{0.49\textwidth}
      \lstcode[basicstyle=\ttfamily\tiny, firstline=19, firstnumber=19]{codeoutput/TreeNode.java}{Java}
    \end{column}
  \end{columns}

  \small \texttt{add} returns \texttt{this}, so calls can be chained. The visitor also gets the \textbf{depth}, for indentation.
\end{frame}
% ~~~~~~~~~~~~~~~~~~~~~~~~~~~~~~~~~~~~~~ A frm




% ~~~~~~~~~~~~~~~~~~~~~~~~~~~~~~~~~~~~~~ V frm
\begin{frame}[t,fragile] \frametitle{A table of contents is an N-ary tree}
  \lstcoderightstep[0.6][basicstyle=\ttfamily\tiny][basicstyle=\ttfamily\tiny]{codeoutput/UseNary.java}{Java}{codeoutput/UseNary.txt}

  \vspace{0.2em}
  \small Preorder with indentation is exactly how a table of contents is printed.
\end{frame}
% ~~~~~~~~~~~~~~~~~~~~~~~~~~~~~~~~~~~~~~ A frm

% ~~~~~~~~~~~~~~~~~~~~~~~~~~~~~~~~~~~~~~ V
\hSection[Summary]{Summary}
% ~~~~~~~~~~~~~~~~~~~~~~~~~~~~~~~~~~~~~~ A




% ~~~~~~~~~~~~~~~~~~~~~~~~~~~~~~~~~~~~~~ V frm
\begin{frame}[t,fragile] \frametitle{Recipes}
  \begin{center}\footnotesize
  \begin{tabular}{ll}
    \toprule
    \textbf{Goal} & \textbf{How} \\
    \midrule
    a tree of any type             & \texttt{class Node<T>}, generic methods \texttt{static <T> \dots} \\
    one traversal, many actions    & pass a \texttt{Consumer<T>} (lambda or method reference) \\
    no recursion                   & preorder/inorder with an \texttt{ArrayDeque} as a stack; levels with a queue \\
    one level per line             & \texttt{for (int k = q.size(); k > 0; k--)} \\
    for-each over a tree           & an \texttt{Iterator} with a stack: $O(h)$ memory \\
    search tree operations         & follow one path: $O(h)$; delete with two children: inorder successor \\
    check a search tree            & pass the range $(lo, hi)$ down \\
    rebuild / save a tree          & preorder + inorder; or preorder with \texttt{null} markers \\
    \bottomrule
  \end{tabular}
  \end{center}
\end{frame}
% ~~~~~~~~~~~~~~~~~~~~~~~~~~~~~~~~~~~~~~ A frm




% ~~~~~~~~~~~~~~~~~~~~~~~~~~~~~~~~~~~~~~ V frm
\begin{frame}[t,fragile] \frametitle{Exercises}
  \begin{enumerate}\small
    \item Write iterative \textbf{postorder} with two stacks (or one stack and a ``last visited'' node).
    \item Add \texttt{min()}, \texttt{max()}, \texttt{floor(k)} and \texttt{height()} to \texttt{Bst}.
    \item Make \texttt{InorderIterator} support \texttt{remove()}. What is difficult?
    \item Rebuild a tree from \textbf{inorder and postorder}. Why do the values have to be distinct?
    \item Write \texttt{TreeNode.levelOrder()} and \texttt{TreeNode.postorder(\dots)}.
    \item Write JUnit tests for \texttt{Bst}: empty tree, one key, duplicates, delete in all three cases.
  \end{enumerate}
\end{frame}
% ~~~~~~~~~~~~~~~~~~~~~~~~~~~~~~~~~~~~~~ A frm




% ~~~~~~~~~~~~~~~~~~~~~~~~~~~~~~~~~~~~~~ V frm
\begin{frame}[t,fragile]
	\begin{center}
		\vspace{4cm}
		\hRemark{\Huge Questions?}
	\end{center}
\end{frame}
% ~~~~~~~~~~~~~~~~~~~~~~~~~~~~~~~~~~~~~~ A frm



\end{document}
